%%%%%%%%%%%%%%%%%%%%%%%%%%%%%%%%%%%%%%%%%
% Beamer Presentation
% LaTeX Template
% Version 1.0 (10/11/12)
%
% This template has been downloaded from:
% http://www.LaTeXTemplates.com
%
% License:
% CC BY-NC-SA 3.0 (http://creativecommons.org/licenses/by-nc-sa/3.0/)
%
%%%%%%%%%%%%%%%%%%%%%%%%%%%%%%%%%%%%%%%%%
%----------------------------------------------------------------------------------------
%	PACKAGES AND THEMES
%----------------------------------------------------------------------------------------

\documentclass{beamer}
\mode<presentation> {
\usetheme{Singapore}
\usecolortheme{rose}
\setbeamertemplate{footline}[page number]
\setbeamertemplate{navigation symbols}{}
}
\usepackage{graphicx} % Allows including images
\usepackage{booktabs} % Allows the use of \toprule, \midrule and \bottomrule in tables
\usepackage{listings}
\usepackage{xcolor}
\usepackage{color}
\colorlet{punct}{red!60!black}
\definecolor{background}{HTML}{EEEEEE}
\definecolor{delim}{RGB}{25,134,57}
\colorlet{numb}{magenta!60!black}
\lstdefinelanguage{json}{
    basicstyle=\normalfont\ttfamily,
    numbers=left,
    numberstyle=\scriptsize,
    stepnumber=1,
    numbersep=8pt,
    showstringspaces=false,
    breaklines=true,
    frame=lines,
    backgroundcolor=\color{background},
    literate=
     *{0}{{{\color{numb}0}}}{1}
      {1}{{{\color{numb}1}}}{1}
      {2}{{{\color{numb}2}}}{1}
      {3}{{{\color{numb}3}}}{1}
      {4}{{{\color{numb}4}}}{1}
      {5}{{{\color{numb}5}}}{1}
      {6}{{{\color{numb}6}}}{1}
      {7}{{{\color{numb}7}}}{1}
      {8}{{{\color{numb}8}}}{1}
      {9}{{{\color{numb}9}}}{1}
      {:}{{{\color{punct}{:}}}}{1}
      {,}{{{\color{punct}{,}}}}{1}
      {\{}{{{\color{delim}{\{}}}}{1}
      {\}}{{{\color{delim}{\}}}}}{1}
      {[}{{{\color{delim}{[}}}}{1}
      {]}{{{\color{delim}{]}}}}{1},
}
%----------------------------------------------------------------------------------------
%	TITLE PAGE
%----------------------------------------------------------------------------------------

\title[AureliaJS]{AureliaJS}
\author{Jesse Javier Cogollo Alvarez}
\institute[EAFIT - TalosDigital]
{
Developer by passion \\
\medskip
\textit{email: cogollo87@gmail.com} \\~\\
\textit{MedellinJS}
}
%\date{\today} % Date, can be changed to a custom date
\begin{document}

\begin{frame}
\titlepage % Print the title page as the first slide
\end{frame}

\begin{frame}
\frametitle{Contenido} % Table of contents slide, comment this block out to remove it
\tableofcontents % Throughout your presentation, if you choose to use \section{} and \subsection{} commands, these will automatically be printed on this slide as an overview of your presentation
\end{frame}

%----------------------------------------------------------------------------------------
%	PRESENTATION SLIDES
%----------------------------------------------------------------------------------------

%------------------------------------------------
\section{AureliaJS} % Sections can be created in order to organize your presentation into discrete blocks, all sections and subsections are automatically printed in the table of contents as an overview of the talk
%------------------------------------------------
\begin{frame}
\frametitle{Que es Aurelia?}
Aurelia es un Framework javascript al lado del servidor. Que maneja simples convenciones.
{\color{blue}\url{http://aurelia.io/}}
\\~\\
Rob Eisenberg
{\color{blue}\url{robeisenberg.com}}
\end{frame}
%------------------------------------------------
\begin{frame}
\frametitle{Caracteristicas}
\begin{columns}[c]
\column{.45\textwidth}
\begin{enumerate}
\item \textbf{Forward-thinking}
\item[•]
\item[•]
\item[•]
\item[•]
\item[•]
\item[•]
\item[•]
\end{enumerate}
\column{.5\textwidth}
Escrito con ES6 y ES7. se integra con Web components. no tiene dependencias externas (excepto polyfills).
\end{columns}
\end{frame}
%------------------------------------------------
\begin{frame}
\frametitle{Caracteristicas}
\begin{columns}[c]
\column{.45\textwidth}
\begin{enumerate}
\item \textbf{Forward-thinking}
\item \textbf{Modern Architecture}
\item[•]
\item[•]
\item[•]
\item[•]
\item[•]
\item[•]
\end{enumerate}
\column{.5\textwidth} % Right column and width
Aurelia esta compuesto por peque\~nos m\'odulos. Utilizados juntos como un completo Framework. Pero tambi\'en permite personalizarlo con diferentes m\'odulos.
\end{columns}
\end{frame}
%------------------------------------------------
\begin{frame}
\frametitle{Caracteristicas}
\begin{columns}[c]
\column{.45\textwidth}
\begin{enumerate}
\item \textbf{Forward-thinking}
\item \textbf{Modern Architecture}
\item \textbf{Two-Way Databinding}
\item[•]
\item[•]
\item[•]
\item[•]
\item[•]
\end{enumerate}
\column{.5\textwidth}
Aurelia tiene una poderosa two-way binding. el cual utiliza t\'ecnicas adaptativos para utilizar el mecanismo mas eficiente para observar cada propiedad. Object.observer
\end{columns}
\end{frame}
%------------------------------------------------
\begin{frame}
\frametitle{Caracteristicas}
\begin{columns}[c]
\column{.45\textwidth}
\begin{enumerate}
\item \textbf{Forward-thinking}
\item \textbf{Modern Architecture}
\item \textbf{Two-Way Databinding}
\item \textbf{Extensible HTML}
\item[•]
\item[•]
\item[•]
\item[•]
\end{enumerate}
\column{.5\textwidth}
La extensibilidad HTML de Aurelia permite crear elementos HTML personalizados, a\~nadir nuevos comportamientos para elementos ya existentes.
\end{columns}
\end{frame}
%------------------------------------------------
\begin{frame}
\frametitle{Caracteristicas}
\begin{columns}[c]
\column{.45\textwidth}
\begin{enumerate}
\item \textbf{Forward-thinking}
\item \textbf{Modern Architecture}
\item \textbf{Two-Way Databinding}
\item \textbf{Extensible HTML}
\item \textbf{Routing and UI Composition}
\item[•]
\item[•]
\item[•]
\end{enumerate}
\column{.5\textwidth}
Router con acoplamiento activo, patrones de ruta din\'amica, activaci\'on de pantallas asincr\'onicamente, etc.
\end{columns}
\end{frame}
%------------------------------------------------
\begin{frame}
\frametitle{Caracteristicas}
\begin{columns}[c]
\column{.45\textwidth}
\begin{enumerate}
\item \textbf{Forward-thinking}
\item \textbf{Modern Architecture}
\item \textbf{Two-Way Databinding}
\item \textbf{Extensible HTML}
\item \textbf{Routing and UI Composition}
\item \textbf{MV* with Conventions}
\item[•]
\item[•]
\end{enumerate}
\column{.5\textwidth} % Right column and width
Quieres invertir tiempo en escribir c\'odigo de configuraci\'on? Aurelia maneja una convenci\'on simple.
\end{columns}
\end{frame}
%------------------------------------------------
\begin{frame}
\frametitle{Caracteristicas}
\begin{columns}[c]
\column{.45\textwidth}
\begin{enumerate}
\item \textbf{Forward-thinking}
\item \textbf{Modern Architecture}
\item \textbf{Two-Way Databinding}
\item \textbf{Extensible HTML}
\item \textbf{Routing and UI Composition}
\item \textbf{MV* with Conventions}
\item \textbf{Broad Language Support}
\item[•]
\end{enumerate}
\column{.5\textwidth}
Utiliza ES5, ES6, TypeScript, AtScript, CoffeeScript.
\end{columns}
\end{frame}
%------------------------------------------------
\begin{frame}
\frametitle{Caracteristicas}
\begin{columns}[c]
\column{.45\textwidth}
\begin{enumerate}
\item \textbf{Forward-thinking}
\item \textbf{Modern Architecture}
\item \textbf{Two-Way Databinding}
\item \textbf{Extensible HTML}
\item \textbf{Routing and UI Composition}
\item \textbf{MV* with Conventions}
\item \textbf{Broad Language Support}
\item \textbf{Testable}
\end{enumerate}
\column{.5\textwidth} % Right column and width
combinando m\'odulos ES6 con contenedores inyecci\'on de dependencias. se hace mas f\'acil crear c\'odigo altamente cohesivo y bajo acoplamiento. haciendo las pruebas unitarias algo f\'acil.
\end{columns}
\end{frame}
%------------------------------------------------
\begin{frame}
\frametitle{Como empezar - Estructura}
Descargamos la ultima versi\'on del skeleton de Aurelia.
\\~\\
{\color{blue}\url{https://github.com/aurelia/skeleton-navigation/releases/tag/0.15.1}}
\\~\\
descomprimimos y ya tenemos una estructura lista para utilizar:
\begin{itemize}
\item build (configuracion gulp).
\item doc (documentacion).
\item src (codificacion).
\item styles (estilos CSS).
\item test (pruebas).
\item README.md (instrucciones para empezar con AureliaJS).
\end{itemize}
\end{frame}
%------------------------------------------------
\begin{frame}
\frametitle{Modulos}
\begin{columns}[c]
\column{.45\textwidth}
\begin{itemize}
\item \textbf{logging}
\end{itemize}
\column{.5\textwidth} % Right column and width
Esta librer\'ia es parte de la plataforma Aurelia y contiene un "Appender" para el logging.
\\~\\
Implementaci\'on:
\textbf{welcome.js}
\end{columns}
\end{frame}
%------------------------------------------------
\begin{frame}
\frametitle{Modulos}
\begin{columns}[c]
\column{.45\textwidth}
\begin{itemize}
\item \textbf{Plugins}
\end{itemize}
\column{.5\textwidth}
Aurelia permite crear e implementar plugins que se adaptan al Framework.
implementaci\'on:
\\~\\
\textbf{index.html}
\\~\\
\textbf{animation-main.js}
\end{columns}
\end{frame}
%------------------------------------------------
\begin{frame}
\frametitle{Vistas y vistas modelos}
\begin{columns}[c]
\column{.45\textwidth}
\begin{itemize}
\item \textbf{DI inyeccion de dependencias}
\end{itemize}
\column{.5\textwidth}
Los AureliaElements son creados como clases los cuales son llamados por el Framework usando contenedor de inyecci\'on de dependencias.
\\~\\
implementaci\'on: \textbf{flickr.js}
\end{columns}
\end{frame}
%------------------------------------------------
\begin{frame}
\frametitle{Vistas y vistas modelos}
\begin{columns}[c]
\column{.45\textwidth}
\begin{itemize}
\item \textbf{Plantillas}
\end{itemize}
\column{.5\textwidth}
El motor de plantillas de Aurelia es el encargado de cargar las vistas y sus recursos, compilando su HTML para un optimo desempe\~no y renderizado del UI en la pantalla.
\\~\\
\textless template\textgreater \textless /template\textgreater. Todo lo que este adentro de esas etiquetas ser\'a administrado por Aurelia.
\\~\\
implementaci\'on: \textbf{welcome.html}
\end{columns}
\end{frame}
\section{DEMO}
%------------------------------------------------
\begin{frame}
\frametitle{Databiding}
\begin{columns}[c]
\column{.45\textwidth}
\begin{itemize}
\item \textbf{Que es}
\end{itemize}
\column{.5\textwidth}
El Databinding nos permite enlazar el estado y comportamiento en un objeto(js) y una vista(html).
\\~\\
Cualquier cambio es enlazado y/o sincronizado en una o varias direcciones. Se puede diferenciar en la vista por que el elemento va seguido de un punto en el atributo con el valor (bind, one-way, two-way o one-time).
\\~\\
implementaci\'on: \textbf{welcome.html}
\end{columns}
\end{frame}
%------------------------------------------------
\begin{frame}
\frametitle{Bindings}
\begin{columns}[c]
\column{.45\textwidth}
\begin{itemize}
\item \textbf{Delegate, trigger and call}
\end{itemize}
\column{.5\textwidth}
Los bindings no solamente conectan elementos o atributos. Tambi\'en son utilizados para lanzar comportamientos.
Se recomienda utilizar el delegate por defecto por que es el mas eficiente al utilizar la memoria y la CPU.
\\~\\
Implementaci\'on: \textbf{welcome.js}
\end{columns}
\end{frame}
%------------------------------------------------
\begin{frame}
\frametitle{Bindings}
\begin{columns}[c]
\column{.45\textwidth}
\begin{itemize}
\item \textbf{String interpolation}
\end{itemize}
\column{.5\textwidth}
Sintaxis: \$\{nameVariable\}. es un binding de un solo camino, la cual la salida se convierte en un string.
\\~\\
Implementaci\'on: \textbf{welcome.html}
\end{columns}
\end{frame}
%------------------------------------------------
\begin{frame}
\frametitle{Bindings}
\begin{columns}[c]
\column{.45\textwidth}
\begin{itemize}
\item \textbf{Ref}
\end{itemize}
\column{.5\textwidth}
Es especial para conectar elementos personalizados, utilizando esta t\'ecnica puedes conectar diferentes componentes.
\\~\\
Implementaci\'on: \textbf{form.html}
\end{columns}
\end{frame}
%------------------------------------------------
\begin{frame}
\frametitle{Bindings}
\begin{columns}[c]
\column{.45\textwidth}
\begin{itemize}
\item \textbf{Select elements}
\end{itemize}
\column{.5\textwidth}
"value.bind" es un HTMLSelectElement que tiene un comportamiento especial para soportar seleccion simple y selecci\'on m\'ultiple de valores. normalmente se combina con el elemento repeat. Pero existe otro elemento especial que permite trabajar con objetos "model.bind".
\\~\\
Implementaci\'on: \textbf{form.html}
\end{columns}
\end{frame}
%------------------------------------------------
\begin{frame}
\frametitle{Bindings}
\begin{columns}[c]
\column{.45\textwidth}
\begin{itemize}
\item \textbf{Radios}
\end{itemize}
\column{.5\textwidth}
"checked.bind" es un HTMLInputElement tiene un comportamiento especial para soportar binding no booleano, como strings y objetos.
\\~\\
Implementaci\'on: \textbf{form.html}
\end{columns}
\end{frame}
%------------------------------------------------
\begin{frame}
\frametitle{Bindings}
\begin{columns}[c]
\column{.45\textwidth}
\begin{itemize}
\item \textbf{Checkboxes}
\end{itemize}
\column{.5\textwidth}
\\~\\
Implementaci\'on: \textbf{form.html}
\end{columns}
\end{frame}
%------------------------------------------------
\begin{frame}
\frametitle{Bindings}
\begin{columns}[c]
\column{.45\textwidth}
\begin{itemize}
\item \textbf{InnerHTML}
\end{itemize}
\column{.5\textwidth}
podemos hacer binding de los elementos innerHTMl con el atributo innerhtml.
\\~\\
Implementaci\'on: \textbf{form.html}
\end{columns}
\end{frame}
%------------------------------------------------
\begin{frame}
\frametitle{Bindings}
\begin{columns}[c]
\column{.45\textwidth}
\begin{itemize}
\item \textbf{textContent}
\end{itemize}
\column{.5\textwidth}
Podemos hacer binding de los elementos textContect con el atributo textcontent.
\\~\\
Implementaci\'on: \textbf{form.html}
\end{columns}
\end{frame}
%------------------------------------------------
\begin{frame}
\frametitle{Bindings}
\begin{columns}[c]
\column{.45\textwidth}
\begin{itemize}
\item \textbf{Style}
\end{itemize}
\column{.5\textwidth}
Podemos hacer binding a los string CSS u objetos en un elemento style.
\\~\\
Implementaci\'on: \textbf{form.html}
\end{columns}
\end{frame}
%------------------------------------------------
\begin{frame}
\frametitle{Bindings}
\begin{columns}[c]
\column{.45\textwidth}
\begin{itemize}
\item \textbf{Adaptive binding}
\end{itemize}
\column{.5\textwidth}
{\color{blue}\url{http://blog.durandal.io/2015/04/03/aurelia-adaptive-binding/}}
\\~\\
Implementaci\'on: \textbf{welcome.html}
\end{columns}
\end{frame}
%------------------------------------------------
\begin{frame}
\frametitle{HTML Extensions}
\begin{columns}[c]
\column{.45\textwidth}
\begin{itemize}
\item \textbf{show}
\end{itemize}
\column{.5\textwidth}
\\~\\
Nos permite condicionar si un elemento o elementos ser\'an visibles o no en el DOM.
Implementaci\'on: \textbf{form.html}
\end{columns}
\end{frame}
%------------------------------------------------
\begin{frame}
\frametitle{HTML Extensions}
\begin{columns}[c]
\column{.45\textwidth}
\begin{itemize}
\item \textbf{if}
\end{itemize}
\column{.5\textwidth}
Tiene el mismo comportamiento que el "show". lo que lo hace diferente es que esta etiqueta si remueve el componente del DOM.
\\~\\
Implementaci\'on: \textbf{form.html}
\end{columns}
\end{frame}
%------------------------------------------------
\begin{frame}
\frametitle{HTML Extensions}
\begin{columns}[c]
\column{.45\textwidth}
\begin{itemize}
\item \textbf{repeat}
\end{itemize}
\column{.5\textwidth}
Este atributo permite renderizar un template m\'ultiples veces y trabaja en conjunto con el ".for"
\\~\\
Implementaci\'on: \textbf{form.html}
\end{columns}
\end{frame}
%------------------------------------------------
\begin{frame}
\frametitle{HTML Extensions}
\begin{columns}[c]
\column{.45\textwidth}
\begin{itemize}
\item \textbf{repeat \$parent}
\end{itemize}
\column{.5\textwidth}
Permite acceder a las propiedades de la viewModel.
%\\~\\
%Implementaci\'on: \textbf{TODO Important}
\end{columns}
\end{frame}
%------------------------------------------------
\begin{frame}
\frametitle{HTML Extensions}
\begin{columns}[c]
\column{.45\textwidth}
\begin{itemize}
\item \textbf{repeat \$index}
\end{itemize}
\column{.5\textwidth}
Es el Index del \'item en el array.
%\\~\\
%Implementaci\'on: \textbf{TODO Important}
\end{columns}
\end{frame}
%------------------------------------------------
\begin{frame}
\frametitle{HTML Extensions}
\begin{columns}[c]
\column{.45\textwidth}
\begin{itemize}
\item \textbf{repeat \$first}
\end{itemize}
\column{.5\textwidth}
Es verdadero si es el primer \'item del array.
%\\~\\
%Implementaci\'on: \textbf{TODO Important}
\end{columns}
\end{frame}
%------------------------------------------------
\begin{frame}
\frametitle{HTML Extensions}
\begin{columns}[c]
\column{.45\textwidth}
\begin{itemize}
\item \textbf{repeat \$last}
\end{itemize}
\column{.5\textwidth}
Es verdadero si es el \'ultimo \'item del array.
%\\~\\
%Implementaci\'on: \textbf{TODO Important}
\end{columns}
\end{frame}
%------------------------------------------------
\begin{frame}
\frametitle{HTML Extensions}
\begin{columns}[c]
\column{.45\textwidth}
\begin{itemize}
\item \textbf{repeat \$even}
\end{itemize}
\column{.5\textwidth}
Es verdadero si el \'item es un n\'umero par.
%\\~\\
%Implementaci\'on: \textbf{TODO Important}
\end{columns}
\end{frame}
%------------------------------------------------
\begin{frame}
\frametitle{HTML Extensions}
\begin{columns}[c]
\column{.45\textwidth}
\begin{itemize}
\item \textbf{repeat \$odd}
\end{itemize}
\column{.5\textwidth}
Es verdadero si el \'item es un n\'umero impar.
%\\~\\
%Implementaci\'on: \textbf{TODO Important}
\end{columns}
\end{frame}
%------------------------------------------------
\begin{frame}
\frametitle{Routing}
\begin{columns}[c]
\column{.45\textwidth}
\begin{itemize}
\item \textbf{Descripci\'on}
\end{itemize}
\column{.5\textwidth}
Cuando implementas enrutadores, tienes varias posibilidades de implementarlo, como: navegaci\'on de la App, dashboards y MDI interfaces.
\\~\\
El Routing de Aurelia, el Router que vive en los viewModels y los router-view que viven en las vistas.
\end{columns}
\end{frame}
%------------------------------------------------
\begin{frame}
\frametitle{Routing}
\begin{columns}[c]
\column{.45\textwidth}
\begin{itemize}
\item \textbf{Que podemos hacer}
\end{itemize}
\column{.5\textwidth}
rutas estaticas: \textbf{/contactUs}
\\~\\
rutas parametriables: \textbf{/user/:userId/dasboard}
\\~\\
rutas wilcard: \textbf{file*path}
\end{columns}
\end{frame}
%------------------------------------------------
\begin{frame}
\frametitle{Routing}
\begin{columns}[c]
\column{.45\textwidth}
\begin{itemize}
\item \textbf{Ciclo de vida}
\\~\\
Existen 4 estados:
\end{itemize}
\column{.5\textwidth}
\textbf{canActivate:} permite validar si se puede navegar en el viewModel.
\\~\\
\textbf{activate:} permite implementar l\'ogica antes del que el viewModel sea cargado.
\\~\\
\textbf{canDeactivate:} permite controlar si deseo que en ese momento se pueda abandonar el viewModel.
\\~\\
\textbf{deactivate:} permite implementar l\'ogica una vez el viewModel ya fue cambiado.
\end{columns}
\end{frame}
%------------------------------------------------
\begin{frame}
\frametitle{Routing}
\begin{columns}[c]
\column{.45\textwidth}
\begin{itemize}
\item \textbf{Convenc\'ion de enrutadores}
\end{itemize}
\column{.5\textwidth}
Implementaci\'on: \textbf{app.js}
\end{columns}
\end{frame}
%------------------------------------------------
\begin{frame}
\frametitle{Routing}
\begin{columns}[c]
\column{.45\textwidth}
\begin{itemize}
\item \textbf{Personalizando el enrutador}
\end{itemize}
\column{.5\textwidth}
Aurelia nos permite agregar pasos en el proceso del Router, los que nos permite agregar validaciones de permisos.
\\~\\
Implementaci\'on: \textbf{app.js}
\end{columns}
\end{frame}
%------------------------------------------------
\begin{frame}
\frametitle{Routing}
\begin{columns}[c]
\column{.45\textwidth}
\begin{itemize}
\item \textbf{pushState}
\end{itemize}
\column{.5\textwidth}
si as\'i lo prefieres, puedes remover el hash de la URL.
\\~\\
Implementaci\'on: \textbf{app.js}
\end{columns}
\end{frame}
%------------------------------------------------
\begin{frame}
\frametitle{Routing}
\begin{columns}[c]
\column{.45\textwidth}
\begin{itemize}
\item \textbf{Reutilizando un viewModel existente}
\end{itemize}
\column{.5\textwidth}
Si alguna vez usted quiere implementar la misma vista para diferentes rutas.
%\\~\\
%Implementaci\'on: \textbf{TODO Important}
\end{columns}
\end{frame}
%------------------------------------------------
\begin{frame}
\frametitle{Routing}
\begin{columns}[c]
\column{.45\textwidth}
\begin{itemize}
\item \textbf{Renderizando multiples viewPorts}
\end{itemize}
\column{.5\textwidth}
Si alguna vez necesitas renderizar contenido en mas de un \'area con la misma ruta.
%\\~\\
%Implementaci\'on: \textbf{TODO Important !!!}
\end{columns}
\end{frame}
%------------------------------------------------
\begin{frame}
\frametitle{Routing}
\begin{columns}[c]
\column{.45\textwidth}
\begin{itemize}
\item \textbf{Rutas personalizadas}
\end{itemize}
\column{.5\textwidth}
Aurelia nos permite generar Routes personalizados.
%\\~\\
%Implementaci\'on: \textbf{TODO Important!!!}
\end{columns}
\end{frame}
%------------------------------------------------
\begin{frame}
\frametitle{Extendiendo HTML}
\begin{columns}[c]
\column{.45\textwidth}
\begin{itemize}
\item \textbf{Que es}
\end{itemize}
\column{.5\textwidth}
Aurelia tiene un compilador de templates. Existen dos tipos de extensiones:
\\~\\
\textbf{Atributos personalizados}
\\~\\
\textbf{Elementos personalizados}
\end{columns}
\end{frame}
%------------------------------------------------
\begin{frame}
\frametitle{Extendiendo HTML}
\begin{columns}[c]
\column{.45\textwidth}
\begin{itemize}
\item \textbf{Atributos personalizados}
\end{itemize}
\column{.5\textwidth}
Los atributos personalizados extienden el HTML con funcionalidades a\~nadiendo nuevos comportamientos para elementos
existentes.
%\\~\\
%Implementaci\'on: \textbf{TODO Important}
\end{columns}
\end{frame}
%------------------------------------------------
\begin{frame}
\frametitle{Extendiendo HTML}
\begin{columns}[c]
\column{.45\textwidth}
\begin{itemize}
\item \textbf{Elementos personalizados}
\end{itemize}
\column{.5\textwidth}
Los elementos personalizados crean nuevos tags al HTML.
%\\~\\
%Implementaci\'on: \textbf{say-hello.js (TODO)}
\end{columns}
\end{frame}
%------------------------------------------------
\begin{frame}
\frametitle{Eventos}
\begin{columns}[c]
\column{.45\textwidth}
\begin{itemize}
\item \textbf{Que es}
\end{itemize}
\column{.5\textwidth}
Es una \'util herramienta  para cuando se necesita desacoplar componentes y necesita hablar con otro componente.
\end{columns}
\end{frame}
%------------------------------------------------
\begin{frame}
\frametitle{Eventos}
\begin{columns}[c]
\column{.45\textwidth}
\begin{itemize}
\item \textbf{Eventos DOM}
\end{itemize}
\column{.5\textwidth}
Debe ser utilizado cuando un mensaje especifico es enviado.
\\~\\
{\color{blue}\url{https://developer.mozilla.org/en-US/docs/Web/Guide/Events/Creating_and_triggering_events}}
\end{columns}
\end{frame}
%------------------------------------------------
\begin{frame}
\frametitle{HTTP Client}
\begin{columns}[c]
\column{.45\textwidth}
\begin{itemize}
\item \textbf{Que es}
\end{itemize}
\column{.5\textwidth}
Aurelia incluye su propia interface para comunicarse con el objeto XMLHttpRequest.
\\~\\
Tiene la mayor\'ia de los verbos est\'andares. adem\'as incluye jsonp.
Cuando realizas una petici\'on, este retorna una promesa con el objeto HttpResponseMessage.
\\~\\
Implementaci\'on: \textbf{flickr.js}
\end{columns}
\end{frame}
%------------------------------------------------
%\begin{frame}
%\frametitle{Personalizaci\'on (Customization)}
%\begin{columns}[c]
%\column{.45\textwidth}
%\begin{itemize}
%\item \textbf{Que es}
%\end{itemize}
%\column{.5\textwidth}
%\\~\\
%Implementaci\'on: \textbf{TODO Important}
%\end{columns}
%\end{frame}
%------------------------------------------------
\begin{frame}
\frametitle{DEMO}
\begin{columns}[c]
\column{.45\textwidth}
\textbf{=)}
\column{.5\textwidth}
\end{columns}
\end{frame}
%------------------------------------------------
\begin{frame}
\frametitle{Redes sociales}
\begin{columns}[c]
\column{.45\textwidth}
\begin{enumerate}
\item \textbf{Twitter}
\item[•]
\item[•]
\end{enumerate}
\column{.5\textwidth}
{\color{blue}@jessecogollo}
\end{columns}
\end{frame}
%------------------------------------------------
\begin{frame}
\frametitle{Redes sociales}
\begin{columns}[c]
\column{.45\textwidth}
\begin{enumerate}
\item Twitter
\item \textbf{Facebook}
\item[•]
\end{enumerate}
\column{.5\textwidth}
{\color{blue}/jessecogollo}
\end{columns}
\end{frame}
%------------------------------------------------
\begin{frame}
\frametitle{Redes sociales}
\begin{columns}[c]
\column{.45\textwidth}
\begin{enumerate}
\item Twitter
\item Facebook
\item \textbf{Github}
\end{enumerate}
\column{.5\textwidth}
{\color{blue}\url{https://github.com/jessecogollo/ama}}
\end{columns}
\end{frame}
%------------------------------------------------
\begin{frame}
\frametitle{Donde aprender}
\begin{columns}[c]
\column{.45\textwidth}
\begin{enumerate}
\item \textbf{Pagina web oficial}
\item[•]
\item[•]
\end{enumerate}
\column{.5\textwidth}
{\color{blue}\url{http://aurelia.io/get-started.html}}
\end{columns}
\end{frame}
%------------------------------------------------
\begin{frame}
\frametitle{Donde aprender}
\begin{columns}[c]
\column{.45\textwidth}
\begin{enumerate}
\item Pagina web oficial
\item \textbf{pluralsight}
\item[•]
\end{enumerate}
\column{.5\textwidth}
{\color{blue}\url{http://www.pluralsight.com/courses/building-applications-aurelia}}
\end{columns}
\end{frame}
%------------------------------------------------
\begin{frame}
\frametitle{Donde aprender}
\begin{columns}[c]
\column{.45\textwidth}
\begin{enumerate}
\item Pagina web oficial
\item pluralsight
\item \textbf{youtube}
\end{enumerate}
\column{.5\textwidth}
{\color{blue}\url{https://www.youtube.com/results?search_query=aurelia+js}}
\end{columns}
\end{frame}
%------------------------------------------------
\begin{frame}
\frametitle{Preguntas}
\begin{figure}
\includegraphics[width=0.4\linewidth]{preguntas.png}
\end{figure}
\end{frame}
%------------------------------------------------
\begin{frame}
\Huge{\centerline{Gracias !!! =)}}
\end{frame}
%----------------------------------------------------------------------------------------
\end{document}